% R6 regime diagram: three small graph topologies (cycle relayed, path
% bridge/cut, cycle severed), each annotated with its residual outcome.
% Structural diagram (node/edge topology + text annotation, no numeric
% axes) so plain tikzpicture per CONVENTION.md, not pgfplots. Topology,
% edge styling (solid/dashed black, dotted grey) and all three annotation
% strings match r6_figures.py's create_regime_diagram() verbatim; the
% residual values 1.225/0.000 are the C6/P6 minimal pair verified in
% \S\ref{sec:mechanism} against sheaf_mechanism_proof.py, and the trial
% counts 200/200, 0/200 are [internal, sheaf_harness_v2.py] as cited there.
% [verified: script r6_figures.py]
\begin{figure}[htbp]
\centering
{\small\bfseries R6 regime --- detection requires a cycle AND relayed reports; cut edges silent, severed edges dark [internal, \texttt{sheaf\_harness\_v2.py}]}\\[6pt]
\begin{minipage}{0.32\textwidth}
\centering
{\small\bfseries (a) cycle $C_6$ relayed}\\[2pt]
\begin{tikzpicture}[
  gnode/.style={circle,draw=black,thick,fill=harborblue,minimum size=5.5mm,inner sep=0pt,font=\tiny\bfseries,text=white},
]
\node[gnode] (n0) at (1.0,0) {0};
\node[gnode] (n1) at (0.5,0.866) {1};
\node[gnode] (n2) at (-0.5,0.866) {2};
\node[gnode] (n3) at (-1.0,0) {3};
\node[gnode] (n4) at (-0.5,-0.866) {4};
\node[gnode] (n5) at (0.5,-0.866) {5};
\draw[black,thick,dashed] (n0) -- (n1);
\draw[black,thick] (n1) -- (n2);
\draw[black,thick] (n2) -- (n3);
\draw[black,thick] (n3) -- (n4);
\draw[black,thick] (n4) -- (n5);
\draw[black,thick] (n5) -- (n0);
\end{tikzpicture}
\\[4pt]
{\color{seagreen}\bfseries\footnotesize residual 1.225\\ detected (200/200\\ cohomology-only)}
\end{minipage}%
\hfill
\begin{minipage}{0.32\textwidth}
\centering
{\small\bfseries (b) path $P_6$ bridge}\\[2pt]
\begin{tikzpicture}[
  gnode/.style={circle,draw=black,thick,fill=harborblue,minimum size=5.5mm,inner sep=0pt,font=\tiny\bfseries,text=white},
]
% Node spacing widened from the original 0.4 to 0.75: at minimum size 5.5mm
% (0.55cm diameter), 0.4cm spacing made every circle overlap its neighbor,
% which swallowed the dashed cut-edge's visible gap entirely (it rendered
% indistinguishable from the solid edges). 0.75cm leaves a real gap between
% circles so the dash pattern on the n4--n5 cut edge actually shows.
\node[gnode] (n0) at (-1.875,0) {0};
\node[gnode] (n1) at (-1.125,0) {1};
\node[gnode] (n2) at (-0.375,0) {2};
\node[gnode] (n3) at (0.375,0) {3};
\node[gnode] (n4) at (1.125,0) {4};
\node[gnode] (n5) at (1.875,0) {5};
\draw[black,thick] (n0) -- (n1);
\draw[black,thick] (n1) -- (n2);
\draw[black,thick] (n2) -- (n3);
\draw[black,thick] (n3) -- (n4);
\draw[black,thick,dashed] (n4) -- (n5);
\end{tikzpicture}
\\[4pt]
{\color{shipred}\bfseries\footnotesize cut edge:\\ residual 0.000\\ provably silent}
\end{minipage}%
\hfill
\begin{minipage}{0.32\textwidth}
\centering
{\small\bfseries (c) cycle $C_6$ severed}\\[2pt]
\begin{tikzpicture}[
  gnode/.style={circle,draw=black,thick,fill=harborblue,minimum size=5.5mm,inner sep=0pt,font=\tiny\bfseries,text=white},
]
\node[gnode] (n0) at (1.0,0) {0};
\node[gnode] (n1) at (0.5,0.866) {1};
\node[gnode] (n2) at (-0.5,0.866) {2};
\node[gnode] (n3) at (-1.0,0) {3};
\node[gnode] (n4) at (-0.5,-0.866) {4};
\node[gnode] (n5) at (0.5,-0.866) {5};
\draw[gray,thick,dotted] (n0) -- (n1);
\draw[black,thick] (n1) -- (n2);
\draw[black,thick] (n2) -- (n3);
\draw[black,thick] (n3) -- (n4);
\draw[black,thick] (n4) -- (n5);
\draw[black,thick] (n5) -- (n0);
\end{tikzpicture}
\\[4pt]
{\color{gray}\bfseries\footnotesize severed:\\ provably dark\\ (0/200)}
\end{minipage}
\caption{Regime diagram for the mechanism: (a) relayed cycle edge --- residual $1.225$, detected
($200/200$ cohomology-only); (b) cut edge --- residual $0.000$, provably silent; (c) severed cycle edge --- provably dark
($0/200$). Detection requires the cycle \emph{and} the relayed reports; each panel is one clause of the observability
contract [internal, \texttt{sheaf\_harness\_v2.py}].}
\label{fig:r6reg}
\end{figure}
